\section{Patch-contribution formulas}
\label{app:contribution-formulas}

Fix a patch based at site~$i$, and write
\begin{equation*}
\alpha_i
=
\sum_{R\subseteq N(i)}\delta_i(R)
+
\sum_{\substack{R\subseteq N(i)\\R\neq\varnothing}}\beta_i(R).
\end{equation*}
Thus~$\alpha_i$ is the total rate of outgoing marked interactions at~$i$, and
$V_i=\alpha_i+a_i^\beta(\varnothing)$. For~$\Delta\in[0,\infty]$ and
$z\in[0,1]$, set
\begin{align}
\varphi_i(\Delta)
&=
e^{-\alpha_i\Delta}
+
\delta_i(\varnothing)
\int_0^\Delta e^{-\alpha_i u}\,du,
\label{eq:patch-phi}
\\
\psi_i(\Delta,z)
&=
\delta_i(\varnothing)
\int_0^\Delta e^{a_i^\beta(\varnothing)u}\,du
+
z e^{a_i^\beta(\varnothing)\Delta}.
\label{eq:patch-psi}
\end{align}
At~$\Delta=\infty$, these expressions are interpreted by their limits.

\subsection{End contributions}
\label{app:end-formulas}

Let~$Q$ be an end patch, and put
\begin{equation*}
i=i(Q),
\qquad
S=S(Q),
\qquad
\Delta=e(Q)-s(Q).
\end{equation*}
Then
\begin{equation}
\label{eq:end-contribution-formulas}
C(z,Q)
=
\begin{cases}
\displaystyle
\frac{\psi_i(\Delta,z)}{\varphi_i(\Delta)},
&\mathsf X(Q)=\mathsf I,
\\[1.4em]
\displaystyle
\frac{
\delta_i(S)\sigma_i^\delta(S)
+
\beta_i(S)\sigma_i^\beta(S)\psi_i(\Delta,z)
}{
\delta_i(S)+\beta_i(S)\varphi_i(\Delta)
},
&\mathsf X(Q)=\mathsf O.
\end{cases}
\end{equation}
In particular, every end contribution is affine in~$z$.

\subsection{Full-patch contributions}
\label{app:bulk-formulas}

Let~$P\in\mathcal P$, and write
\begin{equation*}
i=i(P),
\qquad
S=S(P),
\qquad
\Delta=e(P)-s(P).
\end{equation*}
For the four finite boundary types, and for the two possible infinite types,
\begin{equation}
\label{eq:full-contribution-formulas}
C(P)
=
\begin{cases}
\displaystyle
\frac{\psi_i(\Delta,1)}{\varphi_i(\Delta)},
&\mathsf X(P)\mathsf Y(P)\in\{\mathsf{II},\mathsf{IE}\},
\\[1.4em]
e^{V_i\Delta},
&\mathsf X(P)\mathsf Y(P)=\mathsf{IO},
\\[0.8em]
\displaystyle
\frac{
\delta_i(S)\sigma_i^\delta(S)
+
\beta_i(S)\sigma_i^\beta(S)\psi_i(\Delta,1)
}{
\delta_i(S)+\beta_i(S)\varphi_i(\Delta)
},
&\mathsf X(P)\mathsf Y(P)\in\{\mathsf{OI},\mathsf{OE}\},
\\[1.4em]
\sigma_i^\beta(S)e^{V_i\Delta},
&\mathsf X(P)\mathsf Y(P)=\mathsf{OO}.
\end{cases}
\end{equation}
An infinite full patch has type~$\mathsf{IE}$ or~$\mathsf{OE}$, and the
corresponding row is read at~$\Delta=\infty$. If a denominator in
\eqref{eq:end-contribution-formulas} or
\eqref{eq:full-contribution-formulas} vanishes, that labeled patch has zero
probability of occurring; every realized patch has a positive normalizer.

The same formulas can be written directly in the spin-rate coefficients. The
required substitutions are
\begin{equation}
\label{eq:dual-to-spin-contribution-data}
\begin{aligned}
\alpha_i
&=
\sum_{R\subseteq N(i)}\abs{c_i^0(R)}
+
\sum_{\substack{R\subseteq N(i)\\R\neq\varnothing}}
\abs{c_i^0(R)+c_i^1(R)},
\\
a_i^\beta(\varnothing)
&=-c_i^0(\varnothing)-c_i^1(\varnothing),
\\
\delta_i(S)\sigma_i^\delta(S)&=c_i^0(S),
&\delta_i(S)&=\abs{c_i^0(S)},
\\
\beta_i(S)\sigma_i^\beta(S)&=-c_i^0(S)-c_i^1(S),
&\beta_i(S)&=\abs{c_i^0(S)+c_i^1(S)}.
\end{aligned}
\end{equation}
Since~$c_i^0(\varnothing)\geq0$, this also gives
\begin{align}
\varphi_i(\Delta)
&=
e^{-\alpha_i\Delta}
+
c_i^0(\varnothing)
\int_0^\Delta e^{-\alpha_i u}\,du,
\label{eq:patch-phi-spin}
\\
\psi_i(\Delta,z)
&=
c_i^0(\varnothing)
\int_0^\Delta
e^{-(c_i^0(\varnothing)+c_i^1(\varnothing))u}\,du
+
z e^{-(c_i^0(\varnothing)+c_i^1(\varnothing))\Delta}.
\label{eq:patch-psi-spin}
\end{align}

\subsection{Derivation}
\label{app:contribution-derivation}

Suppose first that a patch starts incoming, so~$X_s^P=1$, and its terminal
label is~$\mathsf I$ or~$\mathsf E$. Consistency occurs either when there is no
outgoing mark before the endpoint, or when the first outgoing mark is an
empty-target death. Therefore
\begin{equation*}
\mathbb P_P(\operatorname{Con}(P))
=
e^{-\alpha_i\Delta}
+
\delta_i(\varnothing)
\int_0^\Delta e^{-\alpha_i u}\,du
=
\varphi_i(\Delta).
\end{equation*}
The corresponding unnormalized weighted expectation with terminal variable
$z$ is
\begin{equation*}
\delta_i(\varnothing)
\int_0^\Delta e^{(V_i-\alpha_i)u}\,du
+
z e^{(V_i-\alpha_i)\Delta}
=
\psi_i(\Delta,z).
\end{equation*}
Dividing gives the first row of~\eqref{eq:end-contribution-formulas}, and
setting~$z=1$ gives the~$\mathsf{II}$ and~$\mathsf{IE}$ row of
\eqref{eq:full-contribution-formulas}.

Suppose next that the patch starts outgoing and has terminal label~$\mathsf I$
or~$\mathsf E$. If~$\alpha(P)=\delta$, the source is inactive immediately
after the initial interaction, so consistency is automatic and the local factor
is~$\sigma_i^\delta(S)$. If~$\alpha(P)=\beta$, the source starts active and the
preceding calculation applies. Using the probabilities in
\eqref{eq:patch-initial-kind-law}, the consistency normalizer is
\begin{equation*}
\frac{
\delta_i(S)+\beta_i(S)\varphi_i(\Delta)
}{
\delta_i(S)+\beta_i(S)
},
\end{equation*}
while the unnormalized weighted expectation is
\begin{equation*}
\frac{
\delta_i(S)\sigma_i^\delta(S)
+
\beta_i(S)\sigma_i^\beta(S)\psi_i(\Delta,z)
}{
\delta_i(S)+\beta_i(S)
}.
\end{equation*}
Their ratio gives the second row of~\eqref{eq:end-contribution-formulas} and the
$\mathsf{OI}$, $\mathsf{OE}$ row of~\eqref{eq:full-contribution-formulas}.

Finally, an outgoing terminal boundary requires the source to remain active
throughout the patch. An incoming-initial patch then contributes
$e^{V_i\Delta}$. For an outgoing-initial patch, consistency also forces
$\alpha(P)=\beta$, leaving~$\sigma_i^\beta(S)e^{V_i\Delta}$. These are the
remaining two rows of~\eqref{eq:full-contribution-formulas}. The same
calculations pass to~$\Delta=\infty$ by taking limits.

\subsection{Empty-neighbour relaxation}
\label{app:positivity-algebra}

Put~$r_i=c_i^0(\varnothing)+c_i^1(\varnothing)$. If~$r_i>0$, set
$q_i=c_i^0(\varnothing)/r_i$. Then
\begin{equation}
\label{eq:empty-neighbour-relaxation}
\psi_i(u,z)
=
q_i+(z-q_i)e^{-r_i u},
\qquad
\psi_i(u+v,z)
=
\psi_i\bigl(u,\psi_i(v,z)\bigr).
\end{equation}
If~$r_i=0$, then~$\psi_i(u,z)=z$, and the second identity remains valid. For
later use, let~$Q$ be an end patch based at~$i$, and let~$u\geq e(Q)$. Write
$Q^{\uparrow u}$ for its continuation through time~$u$ with no intervening
successful interaction, rooted at the original end patch~$Q$. Its contribution,
including the conditional probability of this continuation, is
\begin{equation}
\label{eq:end-patch-continuation}
C(z,Q^{\uparrow u})
=
C\bigl(\psi_i(u-e(Q),z),Q\bigr).
\end{equation}
Since~$C(z,Q)$ is affine,
\begin{equation}
\label{eq:end-patch-slope-relaxation}
\partial_z C(z,Q^{\uparrow u})
=
e^{-r_i(u-e(Q))}\partial_z C(z,Q).
\end{equation}
